\chapter*{About this book}
\addcontentsline{toc}{chapter}{About this book}
Each chapter of this book is one calculation carried out in ECCE from start to finish: the question it answers, what ECCE shows, a numbered procedure and the results. The purpose of the book is to show how ECCE is used for real work. Where a chapter's calculation is also in one of ECCE's example sets, the chapter says so, and the finished calculation can be opened and compared at any step.

The procedures use the menus and panels of ECCE \ecceversion; menu paths are written as \menu{Tools > Diffraction...}, buttons and fields as \ctrl{Calculate}, and commands as \cmd{ecce-xrd}. Everything a procedure does in the windows can also be done from the command line, and AGENTS.md in ECCE's documentation lists the commands.


The Teaching part is written for teachers. Each of its chapters takes a classical discovery or a core concept of chemistry, from the shells of the atom to Photograph 51, and shows how a calculation in ECCE makes it visible in a lecture or a demonstration: what to show, in which order, what the numbers mean and where the picture can mislead.
